\ifajar
  \section{Latihan}
\else
  \section{Bab 5: Performance Engineering}
\fi

Seluruh file yang dipakai pada latihan ini tersedia di
\berkas{sources/chapter05/}, dihitung relatif terhadap akar repositori.
Latihan membutuhkan Docker dengan Compose versi 2, Python 3, dan Node.js
untuk mengompilasi file TypeScript. Listing penyiapan masuk ke direktori
itu lebih dahulu, dan seluruh perintah sesudahnya dijalankan dari sana. Bila
membuka terminal baru, masuk kembali dengan \texttt{cd sources/chapter05}
dari akar repositori.

Peringatan yang sama seperti bab sebelumnya berlaku di sini.
\berkas{sources/chapter05/beban.py} dan
\berkas{sources/chapter05/gate-budget.py} sengaja membebani, sehingga
keduanya hanya diarahkan ke aplikasi milik sendiri. Hasil tiap mahasiswa akan
berbeda, karena angkanya bergantung pada perangkat keras, versi PostgreSQL,
dan beban lain yang sedang berjalan di komputer yang sama. Yang dilaporkan
adalah pola perbandingannya beserta penjelasannya, bukan angka persisnya.

Penyiapannya dilakukan sekali, seperti pada Listing \ref{lst:c5-siapkan}.

\begin{lstlisting}[language=bash, caption={Menyiapkan basis data, data uji, dan file statis, memakai \berkas{sources/chapter05/siapkan-data.sh}}, label={lst:c5-siapkan}]
cd sources/chapter05
docker compose up -d
./siapkan-data.sh
python3 -m venv ../.venv
../.venv/bin/pip install "psycopg[binary,pool]" sqlalchemy
npx -p typescript tsc --target es2022 --module es2022 --strict \
  --lib es2022,dom --outDir statis web-vitals.ts katalog.ts
\end{lstlisting}

Aplikasinya dinyalakan di satu terminal dan dibiarkan menyala selama latihan
berlangsung, seperti pada Listing \ref{lst:c5-aplikasi}.

\begin{lstlisting}[language=bash, caption={Menyalakan aplikasi yang diukur, tersedia di \berkas{sources/chapter05/aplikasi.py}}, label={lst:c5-aplikasi}]
source ../.venv/bin/activate
python aplikasi.py

# Keluarannya:
# Melayani di http://127.0.0.1:8010/  (Ctrl+C untuk berhenti)
\end{lstlisting}

\ifajar
  \subsection*{Latihan 1: Menegakkan Budget Performa}
\else
  \subsection*{Prosedur 1: Menegakkan Budget Performa}
\fi

Latihan ini membuktikan bahwa \textit{budget} hanya berguna bila ada yang
menegakkannya, dan bahwa \textit{verdict}-nya dapat dibaca mesin.

\begin{enumerate}
  \item Baca \berkas{sources/chapter05/budget.json}, lalu tulis dugaan
        untuk tiap alamat: lolos atau tidak, dan mengapa. Dugaan ditulis
        lebih dahulu agar tidak menyesuaikan diri dengan hasilnya.
  \item Jalankan \textit{gate}-nya seperti pada Listing \ref{lst:c5-lat1-jalan},
        lalu catat \textit{verdict} tiap baris ke Tabel \ref{tab:c5-lat1-hasil}. Tabel
        \ref{tab:c5-lat1-contoh} memperlihatkan contoh yang sudah terisi.
  \item Catat status keluar prosesnya dengan \texttt{echo \$?} tepat sesudah
        \textit{gate} selesai. Angka itulah yang dibaca \textit{Continuous
        Integration}.
  \item Ubah satu angka \textit{budget} menjadi longgar, misalnya p95 alamat agregat
        menjadi 2000 milidetik, lalu jalankan ulang \textit{gate}-nya dan catat
        status keluarnya.
  \item Jawab tiga pertanyaan berikut. Pertama, alamat mana yang gagal, dan
        apakah kegagalannya berasal dari waktu atau dari ukuran? Kedua,
        mengapa baris Core Web Vitals dinyatakan gagal padahal aplikasinya
        berjalan normal? Ketiga, apa akibatnya bila \textit{gate} dibuat selalu
        keluar dengan status 0?
\end{enumerate}

\begin{lstlisting}[language=bash, caption={Menjalankan \textit{gate} budget dan membaca status keluarnya, tersedia di \berkas{sources/chapter05/gate-budget.py}}, label={lst:c5-lat1-jalan}]
source ../.venv/bin/activate
python gate-budget.py
echo $?

# Keluarannya:
# Gagal: 6 budget terlampaui.
# 1
\end{lstlisting}

\begin{table}[!htb]
  \centering
  \caption{Contoh hasil Latihan 1 yang sudah terisi, waktu dalam milidetik}
  \label{tab:c5-lat1-contoh}
  \begin{tabularx}{\textwidth}{Xrrl}
    \hline
    \textbf{Ukuran} & \textbf{\textit{Budget}} & \textbf{Terukur} & \textbf{\textit{Verdict}} \\
    \hline
    \texttt{/api/katalog} p95 & 60 & 12,59 & lolos \\
    \hline
    \texttt{/api/katalog} ukuran & 8 KB & 0,29 KB & lolos \\
    \hline
    \texttt{/api/terlaris} p95 & 400 & 830,95 & LEWAT \\
    \hline
    \texttt{/api/katalog-penuh} p95 & 60 & 2.262,70 & LEWAT \\
    \hline
    \texttt{/api/katalog-penuh} ukuran & 8 KB & 275,30 KB & LEWAT \\
    \hline
    LCP p75 & 2500 & tanpa data & LEWAT \\
    \hline
  \end{tabularx}
\end{table}

\begin{table}[!htb]
  \centering
  \caption{Lembar isian hasil pengukuran Latihan 1}
  \label{tab:c5-lat1-hasil}
  \begin{tabularx}{\textwidth}{Xrrl}
    \hline
    \textbf{Ukuran} & \textbf{\textit{Budget}} & \textbf{Terukur} & \textbf{\textit{Verdict}} \\
    \hline
    \texttt{/api/katalog} p95 & \dots & \dots & \dots \\
    \hline
    \texttt{/api/katalog} ukuran & \dots & \dots & \dots \\
    \hline
    \texttt{/api/terlaris} p95 & \dots & \dots & \dots \\
    \hline
    \texttt{/api/katalog-penuh} p95 & \dots & \dots & \dots \\
    \hline
    \texttt{/api/katalog-penuh} ukuran & \dots & \dots & \dots \\
    \hline
    Status keluar sebelum dan sesudah \textit{budget} dilonggarkan & \dots & \dots & \dots \\
    \hline
  \end{tabularx}
\end{table}

\ifajar
  \subsection*{Latihan 2: Menemukan Bagian Termahal}
\else
  \subsection*{Prosedur 2: Menemukan Bagian Termahal}
\fi

Latihan ini membuktikan bahwa dua alamat yang sama-sama lambat dapat
membutuhkan perbaikan yang sama sekali berbeda.

\begin{enumerate}
  \item Jalankan profil kedua alamat seperti pada Listing
        \ref{lst:c5-lat2-jalan}.
  \item Untuk tiap alamat, catat tiga fungsi teratas menurut waktu kumulatif
        beserta bagiannya terhadap waktu total, lalu isi Tabel
        \ref{tab:c5-lat2-hasil}. Tabel \ref{tab:c5-lat2-contoh}
        memperlihatkan contoh yang sudah terisi.
  \item Untuk tiap alamat, tulis satu kalimat perbaikan yang menyentuh
        sebabnya, bukan \textit{symptom}-nya.
  \item Jawab tiga pertanyaan berikut. Pertama, mengapa mempercepat \textit{query}
        tidak banyak menolong alamat katalog penuh? Kedua, mengapa
        mempercepat kode Python tidak menolong alamat agregat? Ketiga,
        berapa kali \texttt{baris\_menjadi\_kamus} dipanggil untuk 20
        ulangan, dan dari mana angka itu berasal?
\end{enumerate}

\begin{lstlisting}[language=bash, caption={Memprofil dua alamat, tersedia di \berkas{sources/chapter05/profil-endpoint.py}}, label={lst:c5-lat2-jalan}]
source ../.venv/bin/activate
python profil-endpoint.py katalog-penuh 20
python profil-endpoint.py terlaris 20
\end{lstlisting}

\begin{table}[!htb]
  \centering
  \caption{Contoh hasil Latihan 2 yang sudah terisi, 20 ulangan}
  \label{tab:c5-lat2-contoh}
  \begin{tabularx}{\textwidth}{lXrr}
    \hline
    \textbf{Alamat} & \textbf{Fungsi teratas} & \textbf{Kumulatif} & \textbf{Bagian} \\
    \hline
    katalog penuh & \texttt{ambil\_katalog\_penuh} & 5,10 s & 68\% \\
    \hline
    katalog penuh & \texttt{json.dumps} & 1,99 s & 26\% \\
    \hline
    katalog penuh & \texttt{baris\_menjadi\_kamus} & 1,61 s & 21\% \\
    \hline
    agregat & \texttt{connection.wait} & 3,84 s & 99\% \\
    \hline
  \end{tabularx}
\end{table}

\begin{table}[!htb]
  \centering
  \caption{Lembar isian hasil pengukuran Latihan 2}
  \label{tab:c5-lat2-hasil}
  \begin{tabularx}{\textwidth}{lXrr}
    \hline
    \textbf{Alamat} & \textbf{Fungsi teratas} & \textbf{Kumulatif} & \textbf{Bagian} \\
    \hline
    katalog penuh & \dots & \dots & \dots \\
    \hline
    katalog penuh & \dots & \dots & \dots \\
    \hline
    katalog penuh & \dots & \dots & \dots \\
    \hline
    agregat & \dots & \dots & \dots \\
    \hline
  \end{tabularx}
\end{table}

\ifajar
  \subsection*{Latihan 3: Menemukan Titik Jenuh}
\else
  \subsection*{Prosedur 3: Menemukan Titik Jenuh}
\fi

Latihan ini mencari kapasitas satu alamat, yaitu jumlah permintaan per detik
yang masih dapat dilayani sebelum \textit{latency}-nya meledak.

\begin{enumerate}
  \item Lakukan \textit{warmup} pada alamatnya dengan 30 permintaan pada
        satu \textit{worker},
        agar \textit{cache} basis data hangat sebelum pengukuran dimulai.
  \item Jalankan beban untuk 1, 2, 4, 8, 16, dan 32 \textit{worker} seperti
        pada Listing \ref{lst:c5-lat3-jalan}, lalu isi Tabel
        \ref{tab:c5-lat3-hasil}. Tabel \ref{tab:c5-lat3-contoh}
        memperlihatkan contoh yang sudah terisi.
  \item Tandai baris tempat \textit{throughput} berhenti bertambah. Itulah
        titik jenuhnya.
  \item Periksa hukum Little pada satu baris pilihan: kalikan
        \textit{throughput} dengan p50 dalam detik, lalu bandingkan hasilnya
        dengan jumlah \textit{worker}.
  \item Jawab tiga pertanyaan berikut. Pertama, pada jumlah \textit{worker}
        berapa \textit{throughput} berhenti bertambah, dan apa yang terjadi
        pada p95 sesudah titik itu? Kedua, mengapa menambah \textit{worker}
        di atas titik jenuh tidak menambah kapasitas? Ketiga, perbaikan mana
        yang menaikkan kapasitas, dan mana yang hanya memindahkan antrean?
\end{enumerate}

\begin{lstlisting}[language=bash, caption={Menaikkan konkurensi bertahap, tersedia di \berkas{sources/chapter05/beban.py}}, label={lst:c5-lat3-jalan}]
source ../.venv/bin/activate
python beban.py /api/terlaris 1 30
python beban.py /api/terlaris 1 15
python beban.py /api/terlaris 8 15
python beban.py /api/terlaris 32 15
\end{lstlisting}

\begin{table}[!htb]
  \centering
  \caption{Contoh hasil Latihan 3 yang sudah terisi, waktu dalam milidetik}
  \label{tab:c5-lat3-contoh}
  \begin{tabularx}{\textwidth}{rrrrX}
    \hline
    \textbf{Worker} & \textbf{p50} & \textbf{p95} & \textbf{Throughput} & \textbf{Catatan} \\
    \hline
    1 & 156,83 & 214,61 & 6,0 & Tanpa antrean \\
    \hline
    8 & 514,98 & 835,65 & 13,6 & Sudah jenuh \\
    \hline
    32 & 2.794,87 & 3.945,97 & 11,0 & \textit{Latency} meledak \\
    \hline
  \end{tabularx}
\end{table}

\begin{table}[!htb]
  \centering
  \caption{Lembar isian hasil pengukuran Latihan 3}
  \label{tab:c5-lat3-hasil}
  \begin{tabularx}{\textwidth}{rrrrX}
    \hline
    \textbf{Worker} & \textbf{p50} & \textbf{p95} & \textbf{Throughput} & \textbf{Catatan} \\
    \hline
    1 & \dots & \dots & \dots & \dots \\
    \hline
    2 & \dots & \dots & \dots & \dots \\
    \hline
    4 & \dots & \dots & \dots & \dots \\
    \hline
    8 & \dots & \dots & \dots & \dots \\
    \hline
    16 & \dots & \dots & \dots & \dots \\
    \hline
    32 & \dots & \dots & \dots & \dots \\
    \hline
  \end{tabularx}
\end{table}

\ifajar
  \subsection*{Latihan 4: Mengumpulkan Core Web Vitals Lapangan}
\else
  \subsection*{Prosedur 4: Mengumpulkan Core Web Vitals Lapangan}
\fi

Latihan ini melengkapi angka lab dengan angka lapangan, yaitu Core Web Vitals
dari kunjungan yang sungguhan.

\begin{enumerate}
  \item Pastikan file TypeScript sudah dikompilasi pada langkah penyiapan,
        lalu buka \texttt{http://localhost:8010/} di browser.
  \item Tekan tombol Muat katalog, tunggu daftarnya tampil, lalu tutup tab
        atau pindah ke tab lain. Laporan dikirim saat halaman ditinggalkan,
        bukan saat selesai dimuat.
  \item Periksa isi \berkas{sources/chapter05/vitals.jsonl} seperti pada
        Listing \ref{lst:c5-lat4-jalan}, lalu catat ketiga metriknya ke Tabel
        \ref{tab:c5-lat4-hasil}. Tabel \ref{tab:c5-lat4-contoh}
        memperlihatkan bentuk isian yang diharapkan.
  \item Ulangi kunjungan sebanyak lima kali, lalu jalankan ulang \textit{gate}
        \textit{budget} dan catat perubahan \textit{verdict} ketiga baris Core Web Vitals.
  \item Jawab tiga pertanyaan berikut. Pertama, mengapa laporan dikirim saat
        halaman ditinggalkan? Kedua, mengapa Core Web Vitals dinilai pada
        persentil 75 dari banyak kunjungan, bukan dari satu kunjungan?
        Ketiga, apa yang terjadi pada CLS bila gambar dipasang tanpa atribut
        \texttt{width} dan \texttt{height}?
\end{enumerate}

\begin{lstlisting}[language=bash, caption={Memeriksa laporan Core Web Vitals yang masuk, dihasilkan \berkas{sources/chapter05/web-vitals.ts}}, label={lst:c5-lat4-jalan}]
cat vitals.jsonl
source ../.venv/bin/activate
python gate-budget.py
\end{lstlisting}

\begin{table}[!htb]
  \centering
  \caption{Bentuk isian Latihan 4 yang diharapkan}
  \label{tab:c5-lat4-contoh}
  \begin{tabularx}{\textwidth}{lrrX}
    \hline
    \textbf{Kunjungan} & \textbf{LCP} & \textbf{INP} & \textbf{CLS} \\
    \hline
    1 & 1.840 ms & 24 ms & 0,0021 \\
    \hline
    p75 dari lima kunjungan & \dots & \dots & \dots \\
    \hline
  \end{tabularx}
\end{table}

\begin{table}[!htb]
  \centering
  \caption{Lembar isian hasil pengukuran Latihan 4}
  \label{tab:c5-lat4-hasil}
  \begin{tabularx}{\textwidth}{lrrX}
    \hline
    \textbf{Kunjungan} & \textbf{LCP} & \textbf{INP} & \textbf{CLS} \\
    \hline
    1 & \dots & \dots & \dots \\
    \hline
    2 & \dots & \dots & \dots \\
    \hline
    3 & \dots & \dots & \dots \\
    \hline
    4 & \dots & \dots & \dots \\
    \hline
    5 & \dots & \dots & \dots \\
    \hline
    p75 dari lima kunjungan & \dots & \dots & \dots \\
    \hline
  \end{tabularx}
\end{table}

Setelah seluruh latihan selesai, container dihentikan dengan
\texttt{docker compose down -v}. Opsi \texttt{-v} ikut menghapus volume
datanya, sehingga penyiapan berikutnya dimulai dari keadaan bersih.
